\documentclass[11pt]{article}
\usepackage[margin=1in]{geometry}
\usepackage{amsmath, amssymb, amsthm}
\usepackage{mathtools}
\usepackage{hyperref}

\theoremstyle{plain}
\newtheorem{theorem}{Theorem}
\newtheorem{lemma}[theorem]{Lemma}
\newtheorem{proposition}[theorem]{Proposition}

\theoremstyle{definition}
\newtheorem{definition}[theorem]{Definition}
\newtheorem{example}[theorem]{Example}

\DeclareMathOperator{\im}{im}
\DeclareMathOperator{\rank}{rank}

\title{Linear Maps and Matrices}
\author{math-foundations / 06-linear-maps}
\date{}

\begin{document}
\maketitle

\subsection*{First, an example}
Before any abstract definition, here is a concrete linear map. Consider rotation
of the plane $\mathbb{R}^2$ counter-clockwise by $90^\circ$. In the standard
basis its matrix is
\[
R_{90} = \begin{pmatrix} 0 & -1 \\ 1 & 0 \end{pmatrix}.
\]
\emph{Read aloud:} ``R-sub-ninety equals the two-by-two matrix with entries zero, minus one, one, zero.''
Apply it to $(1, 0)^\top$:
\[
R_{90} \begin{pmatrix} 1 \\ 0 \end{pmatrix}
= \begin{pmatrix} 0 \\ 1 \end{pmatrix}.
\]
\emph{Read aloud:} ``R-sub-ninety applied to the column vector one, zero equals the column vector zero, one.''
The east-pointing unit vector becomes the north-pointing unit vector. To verify
linearity on $(2, 3)$, note $(2,3) = 2(1,0) + 3(0,1)$, and
\[
R_{90}\begin{pmatrix} 2 \\ 3 \end{pmatrix}
= 2 R_{90}\begin{pmatrix} 1 \\ 0 \end{pmatrix} + 3 R_{90}\begin{pmatrix} 0 \\ 1 \end{pmatrix}
= 2\begin{pmatrix} 0 \\ 1 \end{pmatrix} + 3\begin{pmatrix} -1 \\ 0 \end{pmatrix}
= \begin{pmatrix} -3 \\ 2 \end{pmatrix},
\]
\emph{Read aloud:} ``R-sub-ninety applied to two, three equals two times zero, one plus three times minus-one, zero, which equals minus-three, two.''
matching the direct computation. That property---scalars and sums passing
through the map---is what ``linear'' will mean.

\section{From Vector Spaces to Maps Between Them}

Vector spaces are the stage; linear maps are the actors. A linear map is the
minimal kind of function between two vector spaces that respects both
operations the spaces carry. Once we agree to study only such maps, an
enormous amount of structure becomes available: composition of maps becomes
matrix multiplication, invertibility becomes a determinant condition, and the
problem of solving a linear system becomes the problem of inverting (or
generalised-inverting) a linear map.

Throughout this lesson $\mathbb{F}$ denotes a field (think $\mathbb{R}$ or
$\mathbb{C}$), and all vector spaces are over $\mathbb{F}$.

\begin{definition}[Linear map]
Let $V, W$ be vector spaces over $\mathbb{F}$. A function $T \colon V \to W$
is a \emph{linear map} (or \emph{linear transformation}) if for all
$u, v \in V$ and $\alpha \in \mathbb{F}$,
\[
T(u + v) = T(u) + T(v), \qquad T(\alpha v) = \alpha\, T(v).
\]
\emph{Read aloud:} ``T of u plus v equals T of u plus T of v; and T of alpha times v equals alpha times T of v.''
Equivalently, $T(\alpha u + \beta v) = \alpha T(u) + \beta T(v)$ for all
$\alpha, \beta \in \mathbb{F}$ and $u, v \in V$. The set of linear maps
$V \to W$ is denoted $\mathcal{L}(V, W)$; it is itself a vector space under
pointwise operations.
\end{definition}

\begin{definition}[Kernel and image]
For a linear map $T \colon V \to W$, the \emph{kernel} (or \emph{null
space}) is
\[
\ker T := \{ v \in V : T(v) = 0_W \},
\]
\emph{Read aloud:} ``Kernel of T is defined as the set of v in V such that T of v equals the zero vector of W.''
and the \emph{image} (or \emph{range}) is
\[
\im T := \{ T(v) : v \in V \} \subseteq W.
\]
\emph{Read aloud:} ``Image of T is defined as the set of T of v for v in V, a subset of W.''
Both are subspaces (of $V$ and $W$ respectively). When $V$ is
finite-dimensional, $\rank T := \dim \im T$ and the \emph{nullity} is
$\dim \ker T$.
\end{definition}

\begin{definition}[Matrix representation]
Suppose $V$ has ordered basis $\mathcal{B} = (v_1, \dots, v_n)$ and $W$ has
ordered basis $\mathcal{C} = (w_1, \dots, w_m)$. For each $j$, write
\[
T(v_j) = \sum_{i=1}^m a_{ij}\, w_i,
\]
\emph{Read aloud:} ``T of v-sub-j equals the sum from i equals one to m of a-sub-i-j times w-sub-i.''
which is possible and uniquely so because $\mathcal{C}$ is a basis. The
matrix $[T]_{\mathcal{B}}^{\mathcal{C}} := (a_{ij}) \in \mathbb{F}^{m\times n}$
is the \emph{matrix of $T$ relative to $\mathcal{B}$ and $\mathcal{C}$}.
Writing $[v]_{\mathcal{B}}$ for the coordinate column of $v$ in $\mathcal{B}$,
the defining identity is
\[
[T(v)]_{\mathcal{C}} = [T]_{\mathcal{B}}^{\mathcal{C}} \, [v]_{\mathcal{B}}.
\]
\emph{Read aloud:} ``The coordinate vector of T of v in basis script-C equals the matrix of T from basis script-B to script-C, times the coordinate vector of v in basis script-B.''
\end{definition}

\section{Worked Example}

\begin{example}[Rotation in the plane]
Let $T = R_\theta \colon \mathbb{R}^2 \to \mathbb{R}^2$ rotate vectors
counter-clockwise by $\theta$. In the standard basis $e_1, e_2$,
$T(e_1) = (\cos\theta, \sin\theta)$ and $T(e_2) = (-\sin\theta, \cos\theta)$,
so
\[
[R_\theta] = \begin{pmatrix} \cos\theta & -\sin\theta \\
\sin\theta & \cos\theta \end{pmatrix}.
\]
\emph{Read aloud:} ``The matrix of R-sub-theta equals the two-by-two matrix with entries cosine theta, minus sine theta, sine theta, cosine theta.''
Then $\ker R_\theta = \{0\}$, $\im R_\theta = \mathbb{R}^2$, and the identity
$[R_\alpha][R_\beta] = [R_{\alpha+\beta}]$ recovers the sine and cosine
addition formulas.
\end{example}

\section{The Rank--Nullity Theorem}

\begin{theorem}[Rank--Nullity]\label{thm:rank-nullity}
Let $V$ be a finite-dimensional vector space over $\mathbb{F}$ and
$T \colon V \to W$ a linear map. Then
\[
\dim \ker T + \dim \im T = \dim V.
\]
\emph{Read aloud:} ``The dimension of the kernel of T plus the dimension of the image of T equals the dimension of V.''
\end{theorem}

\begin{proof}
Set $n := \dim V$. Since $\ker T$ is a subspace of the finite-dimensional
space $V$, it is finite-dimensional. Pick a basis $u_1, \dots, u_k$ of
$\ker T$, where $k = \dim \ker T$. By the basis-extension theorem (a
standard consequence of the existence of bases), we can extend this list to
a basis
\[
u_1, \dots, u_k, v_1, \dots, v_r
\]
\emph{Read aloud:} ``u-sub-one through u-sub-k, then v-sub-one through v-sub-r.''
of $V$, with $k + r = n$. We claim $T(v_1), \dots, T(v_r)$ is a basis of
$\im T$.

\emph{Spanning.} Let $w \in \im T$, so $w = T(v)$ for some $v \in V$.
Expand $v = \sum_{i=1}^k \alpha_i u_i + \sum_{j=1}^r \beta_j v_j$. Applying
$T$ and using $T(u_i) = 0$ gives $w = \sum_j \beta_j T(v_j)$, so the
$T(v_j)$ span $\im T$.

\emph{Linear independence.} Suppose
$\sum_{j=1}^r \gamma_j T(v_j) = 0$. By linearity,
$T\!\left( \sum_j \gamma_j v_j \right) = 0$, so
$\sum_j \gamma_j v_j \in \ker T$. Hence there exist
$\delta_1, \dots, \delta_k$ with
$\sum_j \gamma_j v_j = \sum_i \delta_i u_i$, i.e.
$\sum_j \gamma_j v_j - \sum_i \delta_i u_i = 0$. Since
$u_1, \dots, u_k, v_1, \dots, v_r$ is a basis, all coefficients vanish; in
particular every $\gamma_j = 0$.

Therefore $\dim \im T = r$, and
$\dim \ker T + \dim \im T = k + r = n = \dim V$.
\end{proof}

\begin{proposition}[Composition equals matrix product]
Let $S \colon U \to V$ and $T \colon V \to W$ be linear, with chosen bases.
Then $T \circ S$ is linear and
$[T \circ S] = [T] \, [S]$, where the right-hand side is the usual matrix
product.
\end{proposition}

\begin{proof}[Sketch]
Linearity of $T \circ S$ is immediate. Apply both sides to a basis vector of
$U$ and expand using the defining identity
$[T(v)]_{\mathcal{C}} = [T]\,[v]_{\mathcal{B}}$ twice; the matrix-product
formula falls out of the double summation $\sum_k a_{ik} b_{kj}$.
\end{proof}

\section{Consequences for Machine Learning}

The map an attention head computes, $x \mapsto W_O\,(\text{value mix})$, is
a linear map on the residual stream. A rank-$r$ LoRA update $W_O + BA$ with
$A \in \mathbb{R}^{r \times d}$, $B \in \mathbb{R}^{d \times r}$ has update
rank at most $r$; by Theorem~\ref{thm:rank-nullity} applied to $BA$, the
nullity is at least $d - r$. Hence LoRA can only perturb the layer along
an $r$-dimensional subspace of feature directions, the structural fact that
makes it both parameter-efficient and a strong implicit regulariser.

\end{document}
